% Summary of heinz2011: Tier-based strictly local constraints for phonology

\begin{abstract}
\begin{enumerate}
  \item First defines the Tier-based Strictly Local (TSL) class of languages.
  \item Review properties of this class.
  \item Hypothesize that All phonotactic processes are TSL.
\end{enumerate}
\end{abstract}

\section*{motivation}

The TSL class is motivated by tiered phenomena in phonology, which have been studied for a long time and have existing computational formalizations and learning algorithms. However, their formal language properties have not been explored.

\section*{properties of TSL}

\begin{itemize}
  \item TSL contains the Strictly Local class of languages
  \item TSL languages are all Star-free
  \item TSL languages are not contained by the Locally Threshold-Testable or Piecewise Testable classes.
  \item TSL languages are closed under suffix and prefix.
  \item TSL languages are string extension classes, meaning they are:
  \begin{itemize}
    \item closed under intersection
    \item efficiently learnable from positive data (with a known tier)
  \end{itemize}
\end{itemize}

\section*{definition of TSL}

An SL-\textit{k} language L can be defined as (Sigma, S, K) where:

\begin{itemize}
  \item Sigma is a finite alphabet
  \item S is a set of forbidden factors of length \textit{k}, from Sigma*
  \item L contains all strings in Sigma* which don't contain any factor in S
\end{itemize}

A TSL-\textit{k} language L can be defined as (Sigma, T, S, \textit{k}) where:

\begin{itemize}
  \item Sigma is a finite alphabet
  \item T is a subset of Sigma
  \item S is a set of forbidden factors of length \textit{k}, from T*
  \item L contains all strings in Sigma* which, when you remove all symbols not in T from the string, don't contain any factor in S
\end{itemize}

\section*{hypothesis}

All humanly possible segmental phonotactic patterns are TSL.

Three classes of phonotactic constraints:

\begin{enumerate}
  \item Local segmental patterns (SL)
  \item Long-distance segmental patterns (consonant harmony, disharmony, vowel harmony)
  \item Stress patterns
\end{enumerate}

Three kinds of vowel harmony patterns:

\begin{enumerate}
  \item Without neutral vowels
  \begin{enumerate}
    \item analyze same as assimilatory consonant harmony
  \end{enumerate}
  \item With opaque vowels (which don't harmonize, but begin their own harmony domain)
  \begin{enumerate}
    \item analyze by allowing Xi but not iX on the tier, if the opaque vowel is \textit{i}
  \end{enumerate}
  \item With transparent vowels (which don't harmonize or begin their own domain)
  \begin{enumerate}
    \item analyze by removing these from the tier
  \end{enumerate}
\end{enumerate}
